\section{Conclusion}
\label{sec:conclusion}

We asked how much of the cross-project gap in vulnerability detection is
attributable to the program representation rather than the model, and answered
it by specifying an abstraction and measuring it under controls.

$\phi$ is a total labelling from a C grammar's node alphabet onto an operation
taxonomy, with a projection lattice that makes its granularities nested by
construction. Evaluated as a static analysis it exposes failure modes that
downstream accuracy hides; evaluated as a representation, with graph,
backbone, depth, readout and training budget all held fixed, it improves
cross-project ranking by $+0.062$ to $+0.083$ AUC on two corpora and under
three message-passing schemes, winning 22 of 24 project-level folds with every
interval excluding zero, and roughly halves fold-to-fold variance. Under a
fixed representation the three architectures differ by less than half as much
as the representations differ under a fixed architecture, which is the sense
in which we claim the representation matters more.

We are equally clear about what we did not find. A learned edge-weighting
module, an affinity-weighted federated aggregation, and taxonomy refinement
beyond eight types all fail to improve on their ablations. Absolute
performance stays modest --- an AUPRC lift of about $1.9\times$ over base rate
--- and cross-project detection remains open.

The methodological suggestion we would most like to see adopted is the
cheapest one: when a paper introduces a code abstraction, evaluate it as an
analysis. A few hundred annotated call sites caught two defects in our own
mapping that no downstream metric revealed, and that would otherwise have
inflated the statistic used to motivate the design.
